\documentclass[12pt]{article}

\usepackage[T1]{fontenc}
\usepackage[utf8]{inputenc}
\usepackage[brazil]{babel}
\usepackage{lmodern}
\usepackage{microtype}
\usepackage[a4paper,margin=2.5cm]{geometry}
\usepackage{setspace}
\usepackage{indentfirst}

\setstretch{1.15}

\title{Sr. Keynes e os "Clássicos"; Uma Interpretação Sugerida \thanks{Texto Traduzido livremente por Luiz Mario Andrade com o auxílio da ferramenta Chat GPT 5.4}}
\author{J. R. Hicks}
\date{Abril de 1937}

\begin{document}

\maketitle

\begin{center}
\textbf{I}
\end{center}

Será admitido pelo leitor menos caridoso que o valor de entretenimento da \textit{Teoria Geral do Emprego} do Sr. Keynes é consideravelmente realçado pelo seu aspecto satírico. Mas também é claro que muitos leitores ficaram muito perplexos com esta \textit{Dunciad}. Mesmo que estejam convencidos pelos argumentos do Sr. Keynes e reconheçam humildemente terem sido "economistas clássicos" no passado, acham difícil lembrar que acreditavam, em seus dias não regenerados, nas coisas que o Sr. Keynes diz que eles acreditavam. E há, sem dúvida, outros que encontram em suas dúvidas históricas um obstáculo, o que os impede de obter tanta iluminação da teoria positiva quanto poderiam ter obtido de outra forma.

Uma das principais razões para esta situação deve-se, sem dúvida, ao fato de o Sr. Keynes tomar como típico da "economia clássica" os escritos posteriores do Professor Pigou, particularmente \textit{The Theory of Unemployment}. Ora, \textit{The Theory of Unemployment} é um livro bastante novo e um livro extremamente difícil; de modo que se pode dizer com segurança que ainda não causou muita impressão no ensino comum de economia. Para a maioria das pessoas, suas doutrinas parecem tão estranhas e inovadoras quanto as do próprio Sr. Keynes; de modo que ouvir que ele próprio acreditou nessas coisas deixa o economista comum bastante perplexo.

Por exemplo, a teoria do Professor Pigou corre, em uma extensão surpreendente, em termos reais. Não apenas sua teoria é uma teoria de salários reais e desemprego; mas vários problemas que qualquer outra pessoa teria preferido investigar em termos monetários são investigados pelo Professor Pigou em termos de "bens-salário". O economista clássico comum não tem parte neste \textit{tour de force}.

Mas se, em nome do economista clássico comum, declararmos que ele teria preferido investigar muitos desses problemas em termos monetários, o Sr. Keynes responderá que não existe uma teoria clássica de salários nominais e emprego. É bem verdade que tal teoria não pode ser facilmente encontrada nos manuais: mas isso ocorre apenas porque a maioria dos manuais foi escrita em uma época em que as mudanças gerais nos salários nominais em um sistema fechado não representavam um problema importante. Não pode haver dúvida de que a maioria dos economistas pensava ter uma ideia bastante razoável de qual era a relação entre os salários nominais e o emprego.

Nestas circunstâncias, parece valer a pena tentar construir uma teoria "clássica" típica, baseada em um modelo anterior e mais bruto do que o do Professor Pigou. Se pudermos construir tal teoria e mostrar que ela produz resultados que foram, de fato, comumente tomados como certos, mas que não concordam com as conclusões do Sr. Keynes, então teremos finalmente uma base satisfatória de comparação. Podemos esperar ser capazes de isolar as inovações do Sr. Keynes e, assim, descobrir quais são as questões reais em disputa.

Como o nosso propósito é a comparação, tentarei expor a minha teoria clássica típica de forma semelhante àquela em que o Sr. Keynes expõe a sua própria teoria; e deixarei de lado todas as complicações secundárias que não incidam diretamente sobre esta questão especial em mãos. Assim, assumo que estou lidando com um período curto no qual a quantidade de equipamento físico de todos os tipos disponível pode ser considerada fixa. Assumo trabalho homogêneo. Assumo ainda que a depreciação pode ser negligenciada, de modo que a produção de bens de investimento corresponde ao novo investimento. Esta é uma simplificação perigosa, mas as questões importantes levantadas pelo Sr. Keynes em seu capítulo sobre o custo de uso (\textit{user cost}) são irrelevantes para nossos propósitos.

Comecemos assumindo que $w$, a taxa de salários nominais por cabeça, pode ser tomada como dada.

Sejam $x, y$ as produções de bens de investimento e bens de consumo, respectivamente, e $N_x, N_y$ o número de homens empregados na produção dos mesmos. Como a quantidade de equipamento físico especializado para cada indústria é dada, $x = f_x(N_x)$ e $y = f_y(N_y)$, onde $f_x, f_y$ são funções dadas.

Seja $M$ a quantidade dada de moeda.

Deseja-se determinar $N_x$ e $N_y$.

Primeiro, o nível de preços dos bens de investimento = seu custo marginal = $w(dN_x/dx)$. E o nível de preços dos bens de consumo = seu custo marginal = $w(dN_y/dy)$.

A renda auferida nos setores de investimento (valor do investimento, ou simplesmente Investimento) = $wx(dN_x/dx)$. Chamemos isto de $I_x$.

A renda auferida nos setores de consumo = $wy(dN_y/dy)$.

Renda Total = $wx(dN_x/dx) + wy(dN_y/dy)$. Chamemos isto de $I$.

$I_x$ é, portanto, uma função dada de $N_x$, $I$ de $N_x$ e $N_y$. Uma vez que $I$ e $I_x$ sejam determinados, $N_x$ e $N_y$ podem ser determinados.

Agora vamos assumir a "equação da Quantidade de Cambridge" — que existe alguma relação definida entre a Renda e a demanda por moeda. Então, aproximadamente, e à parte o fato de que a demanda por moeda pode depender não apenas da Renda total, mas também de sua distribuição entre pessoas com demandas por saldos relativamente grandes e relativamente pequenas, podemos escrever
\[ M = kI. \]

Assim que $k$ for dado, a Renda total estará determinada.

Para determinar $I_x$, precisamos de duas equações. Uma nos diz que a quantidade de investimento (olhada como demanda por capital) depende da taxa de juros:
\[ I_x = C(i). \]

Isto é o que se torna a escala da eficiência marginal do capital no trabalho do Sr. Keynes.

Além disso, Investimento = Poupança. E a poupança depende da taxa de juros e, se preferir, da Renda. $\therefore I_x = S(i, I)$. (Visto que, no entanto, a Renda já está determinada, não precisamos nos preocupar em inserir a Renda aqui, a menos que queiramos.)

Tomando-os como um sistema, porém, temos três equações fundamentais,
\[ M = kI, \quad I_x = C(i), \quad I_x = S(i, I), \]
para determinar três incógnitas, $I, I_x, i$. Como descobrimos anteriormente, $N_x$ e $N_y$ podem ser determinados a partir de $I$ e $I_x$. O emprego total, $N_x + N_y$, é, portanto, determinado.

Consideremos algumas propriedades deste sistema. Segue-se diretamente da primeira equação que, assim que $k$ e $M$ são dados, $I$ está completamente determinado; ou seja, a renda total depende diretamente da quantidade de moeda. O emprego total, entretanto, não é necessariamente determinado de imediato a partir da renda, uma vez que geralmente dependerá, em certa medida, da proporção da renda poupada e, portanto, da forma como a produção é dividida entre os setores de bens de investimento e de consumo. (Se acontecesse de as elasticidades da oferta serem as mesmas em cada um desses setores, então uma mudança na demanda entre eles produziria movimentos compensatórios em $N_x$ e $N_y$ e, consequentemente, nenhuma mudança no emprego total.)

Um aumento no incentivo ao investimento (ou seja, um movimento para a direita da escala da eficiência marginal do capital, que escrevemos como $C(i)$) tenderá a elevar a taxa de juros e, assim, afetar a poupança. Se o montante da poupança aumentar, o montante do investimento aumentará também; o trabalho será mais empregado nos setores de investimento, menos nos setores de consumo; isto aumentará o emprego total se a elasticidade da oferta nos setores de investimento for maior do que nos setores de bens de consumo — diminuindo-o se for o contrário.

Um aumento na oferta de moeda elevará necessariamente a renda total, pois as pessoas aumentarão seus gastos e empréstimos até que as rendas tenham subido o suficiente para restaurar $k$ ao seu nível anterior. O aumento na renda tenderá a aumentar o emprego, tanto na fabricação de bens de consumo quanto na fabricação de bens de investimento. O efeito total sobre o emprego depende da proporção entre as expansões dessas indústrias; e isso depende da proporção de suas rendas aumentadas que as pessoas desejam poupar, o que também governa a taxa de juros.

Até agora, assumimos que a taxa de salários nominais é dada; mas, enquanto assumirmos que $k$ é independente do nível de salários, não haverá dificuldade com este problema também. Um aumento na taxa de salários nominais diminuirá necessariamente o emprego e elevará os salários reais. Pois uma renda nominal inalterada não pode continuar a comprar uma quantidade inalterada de bens a um nível de preços mais elevado; e, a menos que o nível de preços suba, os preços dos bens não cobrirão seus custos marginais. Deve haver, portanto, uma queda no emprego; à medida que o emprego cai, os custos marginais em termos de trabalho diminuirão e, portanto, os salários reais subirão. (Dado que uma mudança nos salários nominais é sempre acompanhada por uma mudança nos salários reais na mesma direção, se não na mesma proporção, nenhum dano será feito, e alguma vantagem será talvez obtida, se preferirmos trabalhar em termos de salários reais. Naturalmente, a maioria dos "economistas clássicos" seguiu esta linha.)

Creio que se concordará que temos aqui uma teoria razoavelmente consistente, e uma teoria que também é consistente com os pronunciamentos de um grupo reconhecível de economistas. Admitidamente, segue-se desta teoria que se pode ser capaz de aumentar o emprego por meio de inflação direta; mas se você decide ou não favorecer essa política ainda depende do seu julgamento sobre a provável reação sobre os salários e também — em uma área nacional — sobre suas visões sobre o padrão internacional.

Historicamente, esta teoria descende de Ricardo, embora não seja propriamente ricardiana; é provavelmente mais ou menos a teoria que foi defendida por Marshall. Mas com Marshall ela já estava começando a ser qualificada de maneiras importantes; seus sucessores a qualificaram ainda mais. O que o Sr. Keynes fez foi dar uma ênfase enorme às qualificações, de modo que elas quase anulam a teoria original. Sigamos este processo de desenvolvimento.

\begin{center}
\textbf{II}
\end{center}

Quando uma teoria como a teoria "clássica" que acabamos de descrever é aplicada à análise das flutuações industriais, ela encontra dificuldades de várias maneiras. É evidente que a renda nominal total experimenta grandes variações ao longo de um ciclo econômico, e a teoria clássica só pode explicá-las por variações em $M$ ou em $k$, ou, como terceira e última alternativa, por mudanças na distribuição.

(1) A variação em $M$ é a mais simples e óbvia, e tem sido utilizada em grande medida. Mas as variações em $M$ que são rastreáveis durante um ciclo econômico são variações que ocorrem através dos bancos — são variações nos empréstimos bancários; se formos depender delas, é urgentemente necessário que expliquemos a conexão entre a oferta de moeda bancária e a taxa de juros. Isto pode ser feito rudimentarmente pensando nos bancos como pessoas que estão fortemente inclinadas a passar a moeda adiante, emprestando-a em vez de gastá-la. Sua ação, portanto, tende inicialmente a baixar as taxas de juros e, somente depois, quando a moeda passa para as mãos dos gastadores, a elevar os preços e as rendas. "A nova moeda, ou o aumento da moeda, vai não para as pessoas privadas, mas para os centros bancários; e, portanto, aumenta a disposição dos credores em emprestar, em primeira instância, e reduz a taxa de desconto. Mas depois eleva os preços; e, portanto, tende a aumentar o desconto."\footnote{Marshall, \textit{Money, Credit, and Commerce}, p. 257.} Isto é superficialmente satisfatório; mas se tentássemos dar um relato mais preciso deste processo, logo entraríamos em dificuldades. O que determina a quantidade de moeda necessária para produzir uma dada queda na taxa de juros? O que determina a duração do tempo em que a taxa baixa perdurará? Estas não são perguntas fáceis de responder.

(2) Na medida em que dependemos de mudanças em $k$, também podemos ir bem até certo ponto. Mudanças em $k$ podem ser relacionadas a mudanças na confiança, e é realista sustentar que os preços crescentes de um \textit{boom} ocorrem porque o otimismo encoraja uma redução nos saldos; os preços em queda de uma recessão porque o pessimismo e a incerteza ditam um aumento. Mas assim que damos esse passo, torna-se natural perguntar se $k$ não abdicou de seu status de variável independente e não se tornou passível de ser influenciado por outras variáveis em nossas equações fundamentais.

(3) Esta última consideração é poderosamente apoiada por outra, de caráter puramente teórico. Com base na pura teoria do valor, é evidente que o sacrifício direto feito por uma pessoa que detém um estoque de moeda é um sacrifício de juros; e é difícil acreditar que o princípio marginal não opere de todo neste campo. Como Lavington colocou: "A quantidade de recursos que (um indivíduo) detém na forma de moeda será tal que a unidade de moeda que apenas e tão somente vale a pena manter nesta forma lhe rende um retorno de conveniência e segurança igual ao rendimento de satisfação derivado da unidade marginal gasta em bens de consumo, e igual também à taxa líquida de juros."\footnote{Lavington, \textit{English Capital Market}, 1921, p. 30. Ver também Pigou, "The Exchange-value of Legal-tender Money," em \textit{Essays in Applied Economics}, 1922, pp. 179-181.} A demanda por moeda depende da taxa de juros!

O palco está montado para o Sr. Keynes.

\begin{center}
\textbf{III}
\end{center}

Contra as três equações da teoria clássica,
\[ M = kI, \quad I_x = C(i), \quad I_x = S(i, I), \]
o Sr. Keynes começa com três equações,
\[ M = L(i), \quad I_x = C(i), \quad I_x = S(I). \]

Estas diferem das equações clássicas de duas maneiras. Por um lado, a demanda por moeda é concebida como dependendo da taxa de juros (Preferência pela Liquidez). Por outro lado, qualquer possível influência da taxa de juros sobre o montante poupado de uma dada renda é negligenciada. Embora signifique que a terceira equação se torne a equação do multiplicador, que realiza truques tão estranhos, no entanto, esta segunda alteração é uma mera simplificação e, em última análise, insignificante.\footnote{Isto pode ser facilmente visto se considerarmos as equações $M = kI, I_x = C(i), I_x = S(I)$, que incorporam a segunda emenda do Sr. Keynes sem a primeira. A terceira equação já é a equação do multiplicador, mas o multiplicador está com as asas cortadas. Pois, como $I$ ainda depende apenas de $M$, $I_x$ agora depende apenas de $M$, e é impossível aumentar o investimento sem aumentar a disposição de poupar ou a quantidade de moeda. O sistema assim gerado é, portanto, idêntico àquele que, há alguns anos, costumava ser chamado de "Visão do Tesouro". Mas a Preferência pela Liquidez nos transporta da "Visão do Tesouro" para a "Teoria Geral do Emprego".} É a doutrina da preferência pela liquidez que é vital.

Pois agora é a taxa de juros, não a renda, que é determinada pela quantidade de moeda. A taxa de juros confrontada com a escala da eficiência marginal do capital determina o valor do investimento; isso determina a renda pelo multiplicador. Então o volume de emprego (a taxas salariais dadas) é determinado pelo valor do investimento e da renda que não é poupada, mas gasta em bens de consumo.

É este sistema de equações que produz a conclusão surpreendente de que um aumento no incentivo ao investimento, ou na propensão a consumir, não tenderá a elevar a taxa de juros, mas apenas a aumentar o emprego. Apesar disso, porém, e apesar do fato de que uma grande parte do argumento corre em termos deste sistema, e apenas deste sistema, ele não é a \textit{Teoria Geral}. Podemos chamá-lo, se quisermos, de \textit{teoria especial} do Sr. Keynes. A Teoria Geral é algo apreciavelmente mais ortodoxo.

Como Lavington e o Professor Pigou, o Sr. Keynes não acredita, ao fim e ao cabo, que a demanda por moeda possa ser determinada por uma única variável — nem mesmo a taxa de juros. Ele dá mais ênfase a ela do que eles, mas nem para ele nem para eles ela pode ser a única variável a ser considerada. A dependência da demanda por moeda em relação aos juros não faz, no final, mais do que qualificar a antiga dependência em relação à renda. Por mais que enfatizemos o "motivo especulativo", o motivo de "transações" deve sempre entrar também. Consequentemente, temos para a \textit{Teoria Geral}
\[ M = L(I, i), \quad I_x = C(i), \quad I_x = S(I). \]

Com esta revisão, o Sr. Keynes dá um grande passo atrás em direção à ortodoxia marshalliana, e sua teoria torna-se difícil de distinguir das teorias marshallianas revisadas e qualificadas, que, como vimos, não são novas. Existe realmente alguma diferença entre elas, ou é tudo uma luta de fachada? Recorramos a um diagrama (Figura 1).

\begin{center}
\textit{[Espaço para Figura 1: Interseção das curvas IS e LL]}
\end{center}

Contra uma dada quantidade de moeda, a primeira equação, $M = L(I, i)$, nos dá uma relação entre a Renda ($I$) e a taxa de juros ($i$). Esta pode ser desenhada como uma curva ($LL$) que se inclinará para cima, visto que um aumento na renda tende a elevar a demanda por moeda, e um aumento na taxa de juros tende a baixá-la. Além disso, as outras duas equações tomadas em conjunto nos dão outra relação entre Renda e juros. (A escala da eficiência marginal do capital determina o valor do investimento a qualquer taxa de juros dada, e o multiplicador nos diz qual nível de renda será necessário para tornar as poupanças iguais a esse valor do investimento.) A curva $IS$ pode, portanto, ser desenhada mostrando a relação entre Renda e juros que deve ser mantida para tornar a poupança igual ao investimento.

A Renda e a taxa de juros são agora determinadas conjuntamente em $P$, o ponto de interseção das curvas $LL$ e $IS$. Elas são determinadas juntas; assim como o preço e a produção são determinados juntos na teoria moderna da demanda e da oferta. De fato, a inovação do Sr. Keynes é estreitamente paralela, a este respeito, à inovação dos marginalistas. A teoria quantitativa tenta determinar a renda sem os juros, assim como a teoria do valor-trabalho tentou determinar o preço sem a produção; cada uma tem que dar lugar a uma teoria que reconhece um maior grau de interdependência.

Mas se esta é a verdadeira "Teoria Geral", como o Sr. Keynes chega a fazer suas observações sobre um aumento no incentivo ao investimento não elevar a taxa de juros? Pareceria pelo nosso diagrama que um aumento na escala da eficiência marginal do capital deve elevar a curva $IS$; e, portanto, embora eleve a Renda e o emprego, também elevará a taxa de juros.

Isto nos leva ao que, de muitos pontos de vista, é a coisa mais importante no livro do Sr. Keynes. Não é apenas possível mostrar que uma dada oferta de moeda determina uma certa relação entre Renda e juros (que expressamos pela curva $LL$); é também possível dizer algo sobre a forma da curva. Ela provavelmente tenderá a ser quase horizontal à esquerda e quase vertical à direita. Isto ocorre porque há (1) algum mínimo abaixo do qual é improvável que a taxa de juros desça, e (embora o Sr. Keynes não enfatize isso) há (2) um máximo para o nível de renda que pode possivelmente ser financiado com uma dada quantidade de moeda. Se quisermos, podemos pensar na curva como aproximando-se desses limites assintoticamente (Figura 2).

\begin{center}
\textit{[Espaço para Figura 2: Curva LL com parte horizontal e parte vertical]}
\end{center}

Portanto, se a curva $IS$ estiver bem à direita (seja por causa de um forte incentivo ao investimento ou de uma forte propensão a consumir), $P$ recairá sobre aquela parte da curva que é decididamente inclinada para cima, e a teoria clássica será uma boa aproximação, não necessitando de mais do que a qualificação que de fato recebeu nas mãos dos marshallianos posteriores. Um aumento no incentivo ao investimento elevará a taxa de juros, como na teoria clássica, mas também terá algum efeito subsidiário no aumento da renda e, portanto, do emprego também. (O Sr. Keynes em 1936 não é o primeiro economista de Cambridge a ter uma fé temperada em Obras Públicas.) Mas se o ponto $P$ estiver à esquerda da curva $LL$, então a forma especial da teoria do Sr. Keynes torna-se válida. Um aumento na escala da eficiência marginal do capital apenas aumenta o emprego, e não eleva em nada a taxa de juros. Estamos completamente fora de contato com o mundo clássico.

A demonstração deste mínimo é, portanto, de importância central. É tão importante que me aventurarei a parafrasear a prova, expondo-a de uma forma um pouco diferente da adotada pelo Sr. Keynes.\footnote{Keynes, \textit{General Theory}, pp. 201-202.}

Se os custos de manter moeda puderem ser negligenciados, será sempre lucrativo manter moeda em vez de emprestá-la, se a taxa de juros não for superior a zero. Consequentemente, a taxa de juros deve ser sempre positiva. Em um caso extremo, a taxa de curtíssimo prazo pode talvez ser quase zero. Mas, se for assim, a taxa de longo prazo deve situar-se acima dela, pois a taxa longa tem que levar em conta o risco de que a taxa curta possa subir durante a vigência do empréstimo, e deve-se observar que a taxa curta só pode subir, não pode cair.\footnote{É apenas concebível que as pessoas pudessem se tornar tão acostumadas à ideia de taxas curtas muito baixas que não ficariam muito impressionadas com este risco; mas é muito improvável. Pois a taxa curta pode subir, seja porque o comércio melhora e a renda se expande; ou porque o comércio piora e o desejo de liquidez aumenta. Duvido que um sistema monetário tão elástico a ponto de excluir ambas as possibilidades seja realmente concebível.} Isto não significa apenas que a taxa longa deve ser uma espécie de média das taxas curtas prováveis ao longo de sua duração, e que esta média deve situar-se acima da taxa curta atual. Há também o risco mais importante a considerar, de que o credor a longo prazo possa desejar ter dinheiro vivo antes da data acordada para o reembolso e, então, se a taxa curta tiver subido nesse ínterim, ele possa estar envolvido em uma perda substancial de capital. É este último risco que fornece o "motivo especulativo" do Sr. Keynes e que garante que a taxa para empréstimos de duração indefinida (que ele sempre tem em mente como a taxa de juros) não possa cair para muito perto de zero.\footnote{No entanto, algo mais do que o "motivo especulativo" é necessário para explicar o sistema de taxas de juros. A mais curta de todas as taxas curtas deve ser igual à valorização relativa, na margem, da moeda e de tal título; e o título apresenta um desconto principalmente por causa da "conveniência e segurança" de deter moeda — a inconveniência que pode possivelmente ser causada por não ter dinheiro imediatamente disponível. É a chance de que você possa querer descontar o título que importa, não a chance de que você terá que descontá-lo em termos desfavoráveis. O "motivo de precaução", não o "motivo especulativo", é aqui dominante. Mas os termos prospectivos de redesconto são vitais quando se trata da \textit{diferença} entre as taxas curtas e longas.}

Deve-se observar que este mínimo para a taxa de juros se aplica não apenas a uma curva $LL$ (desenhada para corresponder a uma determinada quantidade de moeda), mas a qualquer curva desse tipo. Se a oferta de moeda for aumentada, a curva $LL$ move-se para a direita (como a curva pontilhada na Figura 2), mas as partes horizontais da curva permanecem quase as mesmas. Portanto, novamente, é este marasmo à esquerda do diagrama que perturba a teoria clássica. Se $IS$ estiver à direita, então podemos de fato aumentar o emprego aumentando a quantidade de moeda; mas se $IS$ estiver à esquerda, não podemos fazê-lo; meros meios monetários não forçarão a taxa de juros a baixar mais.

Assim, a Teoria Geral do Emprego é a Economia da Depressão.

\begin{center}
\textbf{IV}
\end{center}

Para elucidar a relação entre o Sr. Keynes e os "Clássicos", inventamos um pequeno aparato. Não parece que tenhamos esgotado os usos desse aparato, por isso concluamos dando-lhe uma pequena corrida por conta própria.

Com esse aparato à nossa disposição, não somos mais obrigados a fazer certas simplificações que o Sr. Keynes faz em sua exposição. Podemos reinserir o $i$ que faltava na terceira equação e levar em conta qualquer efeito possível da taxa de juros sobre a poupança; e, o que é muito mais importante, podemos questionar a dependência exclusiva do investimento em relação à taxa de juros, que parece bastante suspeita na segunda equação. A elegância matemática sugeriria que deveríamos ter $I$ e $i$ em todas as três equações, se a teoria for para ser realmente Geral. Por que não tê-los lá assim:
\[ M = L(I, i), \quad I_x = C(I, i), \quad I_x = S(I, i)? \]

Uma vez que levantamos a questão da Renda na segunda equação, fica claro que ela tem uma ótima pretensão de ser inserida. O Sr. Keynes só consegue deixá-la de fora de forma plausível por meio do seu dispositivo de medir tudo em "unidades de salário", o que significa que ele permite mudanças na escala da eficiência marginal do capital quando há uma mudança no nível de salários nominais, mas que outras mudanças na Renda são consideradas como não afetando a curva, ou pelo menos não da mesma maneira imediata. Mas por que fazer esta distinção? Certamente há todas as razões para supor que um aumento na demanda por bens de consumo, decorrente de um aumento no emprego, frequentemente estimulará diretamente um aumento no investimento, assim que se desenvolva uma expectativa de que o aumento na demanda continuará. Se for assim, deveríamos incluir $I$ na segunda equação, embora deva ser confessado que o efeito de $I$ na eficiência marginal do capital será instável e irregular.

A Teoria Geral Generalizada pode então ser exposta desta forma. Assuma primeiro uma renda nominal total dada. Desenhe uma curva $CC$ mostrando a eficiência marginal do capital (em termos nominais) naquela Renda dada; uma curva $SS$ mostrando a curva de oferta de poupança naquela Renda dada (Figura 3). Sua interseção determinará a taxa de juros que torna as poupanças iguais ao investimento naquele nível de renda. A isto podemos chamar a "taxa de investimento".

Se a Renda subir, a curva $SS$ mover-se-á para a direita; provavelmente $CC$ também se moverá para a direita. Se $SS$ se mover mais do que $CC$, a taxa de juros de investimento cairá; se $CC$ mais do que $SS$, ela subirá. (O quanto ela sobe ou cai, no entanto, depende das elasticidades das curvas $CC$ e $SS$.)

A curva $IS$ (desenhada em um diagrama separado) mostra agora a relação entre a Renda e a correspondente taxa de juros de investimento. Ela tem que ser confrontada (como em nossas construções anteriores) com uma curva $LL$ que mostra a relação entre a Renda e a taxa de juros "monetária"; apenas agora podemos generalizar um pouco a nossa curva $LL$. Em vez de assumir, como antes, que a oferta de moeda é dada, podemos assumir que existe um dado sistema monetário — que até certo ponto, mas apenas até certo ponto, as autoridades monetárias preferirão criar nova moeda em vez de permitir que as taxas de juros subam. Tal curva $LL$ generalizada inclinar-se-á para cima apenas gradualmente — a elasticidade da curva dependendo da elasticidade do sistema monetário (no sentido monetário comum).

\begin{center}
\textit{[Espaço para Figura 3: Curvas CC e SS cruzando-se, e gráfico IS-LL ao lado]}
\end{center}

Como antes, a Renda e os juros são determinados onde as curvas $IS$ e $LL$ se cruzam — onde a taxa de juros de investimento é igual à taxa monetária. Qualquer mudança no incentivo ao investimento ou na propensão a consumir deslocará a curva $IS$; qualquer mudança na preferência pela liquidez ou na política monetária deslocará a curva $LL$. Se, como resultado de tal mudança, a taxa de investimento for elevada acima da taxa monetária, a Renda tenderá a subir; no caso oposto, a Renda tenderá a cair; a extensão em que a Renda sobe ou cai depende das elasticidades das curvas.\footnote{Visto que $C(I, i) = S(I, i)$,
\[ \frac{dI}{di} = - \frac{\partial S / \partial i - \partial C / \partial i}{\partial S / \partial I - \partial C / \partial I} \]
O mercado de poupança e investimento não será estável a menos que $\partial S / \partial i + (- \partial C / \partial i)$ seja positivo. Creio que podemos assumir que esta condição é cumprida. Se $\partial S / \partial i$ for positivo, $\partial C / \partial i$ negativo, $\partial S / \partial I$ e $\partial C / \partial I$ positivos (o estado de coisas mais provável), podemos dizer que a curva $IS$ será mais elástica quanto maiores forem as elasticidades das curvas $CC$ e $SS$, e maior for $\partial C / \partial I$ em relação a $\partial S / \partial I$. Quando $\partial C / \partial I > \partial S / \partial I$, a curva $IS$ é inclinada para cima.}

Quando generalizada desta forma, a teoria do Sr. Keynes começa a parecer muito com a de Wicksell; isso, é claro, não é surpreendente.\footnote{Cf. Keynes, \textit{General Theory}, p. 242.} Há de fato um caso especial onde ela se ajusta perfeitamente à construção de Wicksell. Se houver "pleno emprego" no sentido de que qualquer aumento na Renda provoca imediatamente um aumento nas taxas de salários nominais; então é possível que as curvas $CC$ e $SS$ sejam movidas para a direita exatamente na mesma medida, de modo que $IS$ seja horizontal. (Digo possível porque não é improvável, de fato, que o aumento no nível salarial possa criar uma presunção de que os salários subirão novamente mais tarde; se assim for, $CC$ provavelmente será deslocada mais do que $SS$, de modo que $IS$ será inclinada para cima.) Seja como for, se $IS$ for horizontal, temos uma construção perfeitamente wickselliana;\footnote{Cf. Myrdal, "Gleichgewichtsbegriff," em \textit{Beiträge zur Geldtheorie}, ed. Hayek.} a taxa de investimento torna-se a taxa natural de Wicksell, pois neste caso pode-se pensar nela como determinada por causas reais; se houver um sistema monetário perfeitamente elástico, e a taxa monetária for fixada abaixo da taxa natural, há inflação cumulativa; deflação cumulativa se for fixada acima.

Isto, entretanto, é visto agora como sendo apenas um caso especial; podemos usar nossa construção para abrigar possibilidades muito mais amplas. Se houver muito desemprego, é muito provável que $\partial C / \partial I$ seja bastante pequeno; nesse caso, pode-se confiar que a $IS$ se incline para baixo. Este é o tipo de Economia da Recessão com a qual o Sr. Keynes está amplamente preocupado. Mas não se pode escapar da impressão de que pode haver outras condições em que as expectativas são inflamáveis, quando uma leve tendência inflacionária as acende muito facilmente. Então $\partial C / \partial I$ pode ser grande e um aumento na Renda tender a elevar a taxa de juros de investimento. Nestas circunstâncias, a situação é instável a qualquer taxa monetária dada; é apenas um sistema monetário imperfeitamente elástico — uma curva $LL$ ascendente — que pode impedir a situação de sair de controle completamente.

Estas, então, são algumas das coisas que podemos obter do nosso aparato esquelético. Mas mesmo que ele possa pretender ser uma ligeira extensão do esqueleto semelhante do Sr. Keynes, ele permanece um assunto terrivelmente bruto e improvisado. Em particular, o conceito de "Renda" é trabalhado de forma monstruosa; a maioria de nossas curvas não é realmente determinada a menos que algo seja dito sobre a distribuição da Renda, bem como sobre sua magnitude. De fato, o que elas expressam é algo como uma relação entre o sistema de preços e o sistema de taxas de juros; e não se pode colocar isso em uma curva. Além disso, todos os tipos de questões sobre depreciação foram negligenciados; e todos os tipos de questões sobre o tempo dos processos sob consideração.

A \textit{Teoria Geral do Emprego} é um livro útil; mas não é nem o começo nem o fim da Economia Dinâmica.

\begin{flushright}
J. R. HICKS \\
Gonville and Caius College \\
Cambridge
\end{flushright}

\end{document}